% Compact UCAS report layout, inspired by the course template of
% Shing-Ho Lin and Jun-Xiong Ji (2022). This adaptation is not a new official template.
% Uses installed TeX Fandol fonts; no private/proprietary font copies required.
\documentclass[UTF8,11pt,fontset=fandol]{ctexart}
\usepackage[a4paper,margin=2cm,top=2.2cm,bottom=2.2cm]{geometry}
\usepackage{amsmath,amssymb,graphicx,booktabs,longtable,tabularx,float}
\usepackage[hidelinks]{hyperref}
\setlength{\parindent}{2em}
\setlength{\parskip}{0.25em}
\setlength{\intextsep}{8pt plus 2pt minus 2pt}
\setlength{\emergencystretch}{2em}
\renewcommand{\arraystretch}{1.18}
\newcommand{\Result}[1]{\csname result@#1\endcsname}
\newcommand{\ResultUncertainty}[1]{\csname uncertainty@#1\endcsname}
\input{generated/values.tex}
\begin{document}
\pagestyle{plain}
\begin{center}
{\LARGE\bfseries 《基础物理实验》实验报告}\\[0.8em]
{\large\bfseries @@EXPERIMENT@@}\\[0.5em]
\begin{tabularx}{\textwidth}{@{}X X@{}}
姓名：@@STUDENT_NAME@@ & 学号：@@STUDENT_ID@@\\
分班分组：@@GROUP@@ \quad 座号：@@SEAT@@ & 指导教师：@@TEACHER@@\\
实验日期：@@DATE@@ & 实验地点：@@ROOM@@
\end{tabularx}
\par\smallskip\hrule
\end{center}
@@STATUS@@
\input{content.tex}
\clearpage
\appendix
\input{generated/appendix.tex}
\end{document}
